


The UniZ admin portal gives non-student staff a dedicated interface for managing approvals, academic data, campus security, and announcements. Every administrative role — caretaker, warden, SWO, dean, director, and webmaster — uses the same portal with a role-appropriate view.

\subsection{Logging in}

All admin accounts use the same login endpoint as students. The `role` field in the response determines which dashboard and features you can access.

**Endpoint:** `POST /auth/login`

\begin{lstlisting}[language=bash]
POST https://api.uniz.rguktong.in/api/v1/auth/login
Content-Type: application/json

\end{lstlisting}

\begin{lstlisting}[language=json]
{
  "username": "warden_male",
  "password": "your_password"
}

\end{lstlisting}

**Response:**

\begin{lstlisting}[language=json]
{
  "success": true,
  "token": "eyJ...",
  "role": "warden_male",
  "username": "warden_male"
}

\end{lstlisting}

Store the returned token and include it as a `Bearer` token in the `Authorization` header for all subsequent requests.

<Info>
  Admin usernames follow a fixed pattern: `director`, `dean_cse`, `warden_male`,
  `warden_female`, `caretaker_male`, `caretaker_female`, `security_admin`, and
  `webmaster`. Contact your system administrator if you do not have your
  credentials.
</Info>

\subsection{Role-based access}

Each admin role has a tailored view of the portal. The table below maps roles to their primary responsibilities.

| Role      | Username pattern                      | Primary access                                        |
| --------- | ------------------------------------- | ----------------------------------------------------- |
| Director  | `director`                            | Final approvals, student search, academic publishing  |
| Dean      | `dean_cse`                            | Final approvals, grade management, academic oversight |
| Warden    | `warden_male` / `warden_female`       | Outpass approvals (second level)                      |
| Caretaker | `caretaker_male` / `caretaker_female` | Outpass approvals (first level)                       |
| SWO       | `swo`                                 | Outpass approvals (third level)                       |
| Security  | `security_admin`                      | Gate check-in and check-out                           |
| Webmaster | `webmaster`                           | Banner management, grade uploads                      |

\subsection{What the dashboard shows}

When you log in, the dashboard surfaces the information most relevant to your role.

- **Caretakers and wardens** see a queue of pending outpass requests assigned to their level.
- **Security staff** see the access control terminal with real-time student presence data.
- **Dean and director** see pending final-approval requests, student search, and academic controls.
  - **Webmaster** sees banner management and announcement tools.

The portal automatically hides features your role cannot access. You do not need to configure anything — access is determined by the `role` returned at login.

\subsection{System health}

You can check the status of all UniZ backend services at any time without authentication.

**Endpoint:** `GET /system/health`

\begin{lstlisting}[language=bash]
GET https://api.uniz.rguktong.in/api/v1/system/health

\end{lstlisting}

\begin{lstlisting}[language=json]
{
  "success": true,
  "services": [
    { "name": "Auth Service", "status": "UP", "latency": "120ms" },
    { "name": "Mail Service", "status": "UP", "latency": "45ms" }
  ]
}

\end{lstlisting}

Use this endpoint to verify that the platform is operational before performing bulk operations like grade uploads or result publishing.

\subsection{Key admin sections}


\begin{tcolorbox}[colback=gray!5,colframe=primary,title=Documentation Note]

  
\begin{tcolorbox}[colback=gray!5,colframe=primary,title=Documentation Note]

    Review, approve, forward, or reject student outpass and outing requests.
    Each role operates at a specific level in the multi-step chain.
  
\end{tcolorbox}

  <Card
    title="Academic management"
    icon="graduation-cap"
    href="/admin/academics"
  >
    View batch grades, run bulk updates, upload students, manage semesters, and
    publish results by email.
  
\end{tcolorbox}

  
\begin{tcolorbox}[colback=gray!5,colframe=primary,title=Documentation Note]

    Check students in and out at the campus gate, view approved passes, and
    monitor who is currently outside campus.
  
\end{tcolorbox}

  
\begin{tcolorbox}[colback=gray!5,colframe=primary,title=Documentation Note]

    Create and publish homepage announcements visible to all students in the
    portal.
  
\end{tcolorbox}

</CardGroup>


\section{Deep Technical Analysis}
The underlying system architecture for this module is built upon the \textbf{UniZ Microservices Framework}. 
This design ensures that each functional unit (Auth, Profile, Academics) remains isolated, preventing cascading failures across the university ecosystem.

\subsection{Operational Hardening}
Deployments are managed via K3s (Kubernetes) and containerized using high-performance Alpine Linux Docker images. 
Resource management is critical; for instance, the documentation service itself is provisioned with 2GB of RAM to ensure the responsive rendering of complex documentation graphs and state-management components.

\subsection{Enterprise Security Topology}
Every request entering the UniZ gateway is subjected to an aggressive JWT validation layer. 
Permissions are granular, mapping directly onto the RBAC (Role-Based Access Control) matrix defined in the core system specifications.

\subsection{Data Governance}
Prisma-driven data access ensures type safety and predictable query performance. 
We utilize PostgreSQL connection pooling to handle concurrent requests from the student portal and administrative dashboards simultaneously.
